\documentclass[11pt,a4paper]{article}
\usepackage[margin=25mm]{geometry}
\usepackage{amsmath,amssymb,amsthm,booktabs}
\newtheorem{lemma}{Lemma}\newtheorem{theorem}{Theorem}
\title{The order 14 column from intersections of projective maximals}
\author{Human Compression project}\date{}
\begin{document}\maketitle
Let $G=O'N$. As before, $x_J(p,q)$ counts ordered pairs $(r,s)$
in $J$ of orders $p,q$ with $|rs|=2$, and $e_J(p,q)$ counts those
generating $J$. We extend the accepted four-entry order 14 argument
to the entire column by retaining the flag-normalizer contribution.
The complete subgroup roster, primary ATLAS centralizer data and
character tables remain the same explicit borrowed premises.

We recall the accepted structures and incidences. Each order 14
element $s$ has cyclic centralizer of order 28, with unique involution
$t=s^7$. It lies in a unique row 7 maximal subgroup $C_G(t)$,
and in exactly one projective maximal of each class in which it
is inner. Call these two canonical maximals $M_2(s),M_3(s)$.
Each is isomorphic to $L_3(7):2$, where inner means membership
in the simple socle. The only maximal classes containing an
element of order 14 are rows 2, 3 and 7, as the complete cyclic
order 14 ancestor list shows.

A row 22 flag normalizer $B$ has order 8232 and unique projective
maximal overgroups $P_2,P_3$, of which it is the intersection.
Their uniqueness follows from
$122760\cdot456=55978560$, the class length of row 22.
The accepted inner-involution centralizer is the order 1344
block group $H$; the outer-involution centralizer is
$K=\operatorname{PGL}_2(7)\times C_2$, of order 672.

\begin{lemma}\label{cover}
Every element of order 14 in $B$ is inner in at least one of
$P_2,P_3$. It is inner in both precisely when its cyclic
centralizer of order 28 is contained in $B$.
\end{lemma}
\begin{proof}
The retained complete flag-normalizer model has 1470 elements of
order 14, of which 882 are in the projective socle. This count
applies to either parent: in each projective maximal the row 22
subgroups form the single class of normalizers of transvection
subgroups, whose explicit model was already identified.

There are 588 elements of order 28 in $B$. Their squares form
294 distinct elements of order 14. Indeed the literal reader
checks this image directly; alternatively, any order 14 element
has exactly two order 28 square roots in its cyclic global
centralizer. Thus exactly 294 such elements have their full
cyclic centralizer contained in $B$.

If $s$ is inner in both parents, they are its two canonical
projective maximals, so both contain $C_G(s)$; their intersection
$B$ contains it too. Conversely, if $C_G(s)\leq B$, then $s$ is
a square in both projective maximals and hence is inner in both.
The two sets of inner elements consequently have intersection
of size 294, and their union has size
$882+882-294=1470$. This is the whole order 14 set of $B$.
\end{proof}

\begin{lemma}\label{proper}
Fix a projective maximal $M$ and an outer element $s\in M$ of
order 14. Every proper pair $(r,s)$ with $|rs|=2$ lies either
in $K_s=C_M(s^7)$ or in a row 22 subgroup containing $s$.
There is at most one such row 22 subgroup in $M$.
\end{lemma}
\begin{proof}
The complete maximal list of $M$ has orders
$1876896,8232,1344,672,144,144,114$. The first is its socle,
which excludes outer $s$. An order 1344 inner-involution
centralizer cannot contain $s$, since its central inner involution
would have to equal the unique commuting involution $s^7$,
which is outer. Orders 144 and 114 exclude seven. The remaining
types are exactly the flag normalizer and the outer-involution
centralizer. Every order 672 subgroup containing $s$ is $K_s$,
by the same uniqueness of its central involution.

Suppose a row 22 subgroup $B\leq M$ contains $s$. Its other
projective parent is inner on $s$, by Lemma~\ref{cover}, and
therefore is the unique canonical parent of the other class.
Since $B$ is maximal in $M$, it is the intersection of $M$
with that canonical parent. This determines $B$ uniquely.
Finally $B\cap K_s=C_B(s^7)$, which gives an explicit description
of the only overlap between the two proper contributions.
\end{proof}

Define
\[
y_B(p,14)=|\{(r,s)\in B^2: |r|=p,\ |s|=14,\ |rs|=2,
                                      \ rs^7=s^7r\}|.
\]
\begin{theorem}\label{count}
For every order $p$ in the fixed source table,
\begin{align*}
e_G(p,14)={}&x_G(p,14)-\sum_M[G:M]x_M(p,14)\\
 &+55978560\bigl(x_B(p,14)-y_B(p,14)\bigr)\\
 &+2\cdot342868680x_H(p,14)
   +4\cdot342868680x_K(p,14).
\end{align*}
The sum is over the thirteen maximal conjugacy classes.
\end{theorem}
\begin{proof}
Start with the sum of all maximal contributions. For an outer
$s$ in a projective maximal, remove every pair generating a
proper subgroup of that maximal. By Lemma~\ref{proper}, these
are the union of its unique possible flag-normalizer contribution
and its unique outer-involution centralizer contribution.
They are already covered, respectively, by the canonical
projective parent and by $C_G(s^7)$. All remaining outer pairs
generate their projective maximal, and thus overlap no other
maximal contribution.

There are 5586 outer-involution centralizers in each projective
maximal and 122760 maximals in each of its two classes. Their
full removal contributes
$2\cdot122760\cdot5586x_K=4\cdot342868680x_K$ to the
generating-pair formula. This counts actual parent incidences;
no uniqueness of a global order 672 subgroup's parent is assumed.

After these removals, only the canonical projective maximals
and the unique row 7 maximal can overlap. Their two intersections
with row 7 are the block centralizers in rows 65 and 66. Each
has class length 342868680 and unique actual projective and
row 7 parents, by the accepted direct incidence identities.
Their correction is therefore $2\cdot342868680x_H$.

It remains to combine the flag-normalizer terms. Fix an actual
$B$ and a counted pair $(r,s)\in B^2$. If $s$ is inner in just
one parent, its occurrence in the outer parent is removed once.
If $r$ commutes with $s^7$, it was also counted in that parent's
$K_s$ and must be subtracted once from the removal.
If instead $s$ is inner in both parents, the pair occurs in the
intersection of the two canonical projective maximals, giving
one correction. When $r$ commutes with $s^7$, it is also in
row 7 and is a triple intersection, requiring one subtraction.
Lemma~\ref{cover} excludes every other case. In both cases the
net coefficient is one minus the indicator of $rs^7=s^7r$.
There are 55978560 actual subgroups $B$, proving the asserted
term $55978560(x_B-y_B)$ and the formula.
\end{proof}

The literal reader reuses the full 8232 and 1344 permutation sets
and the complete 672-element product. It checks all order 14
classes, every partner, and the commuting condition defining
$y_B$. The previously retained output already recorded these
conditions; no new group acquisition is required. The nonzero
values of $y_B(p,14)$ occur at $p=2,7$, and are 24696 and 1470,
respectively. For the three previously unpaid entries the data are
\[
\begin{array}{c|rrrr|r}
p&x_B&y_B&x_H&x_K&e_G(p,14)/(2|G|)\\\hline
4&16464&0&0&672&166\\
7&72030&1470&48&384&1054\\
14&148176&0&0&336&1796.
\end{array}
\]
Substitution of the acquired character values gives the entire column:
\[
\begin{array}{c|rrrrrrrr}
p&3&4&5&6&7&8&10&11\\\hline
e_G/(2|G|)&10&166&290&597&1054&3056&2526&4601\\[2pt]
p&12&14&15&16&19&20&28&31\\\hline
e_G/(2|G|)&1391&1796&2330&12504&7899&5024&3616&3266.
\end{array}
\]
The automorphism action on generating pairs is free, and there
is no division by duality. The thirteen previously accepted
entries agree with their retained proofs, and all sixteen agree
with the fixed printed table.

\paragraph{Attribution and scope.}
The numerical assertions are from Connor and Leemans,
\emph{Algorithmic enumeration of regular maps}, Table~1.
We retain the exact source, classification and character premises
of the earlier components. The additional argument is the covering
of the flag normalizer's order 14 elements by its two projective
socles, followed by the uniform incidence correction above.
This component adds three numerical recoveries, without claiming
new chirality, priority for elementary group theory, or completion
of the original manuscript.
\end{document}
